\begin{lemma}[Interchangeability on a closed PE core]
Let $T$ be a positive-rung extended Tashkinov tree constructed from a
$k$-triple $G,e_0,\varphi$, with arbitrary certified input coloring,
terminal ladder tree $T_n$, terminal generating coloring $\varphi_n$,
and terminal stored-color set $D_n$.  Assume that $T$ has maximum
property and that its terminal extension is a PE.  Let $R_n$ and
$T_n\vee R_n$ be the closure and exact ordered union in a hierarchy
base under $\varphi_n$ (respectively `base.closure' and `base.core'
in Lean).

Assume, uniformly over every certified input coloring and generating
coloring, that every positive-rung ETT $T'$ of rung number $s<n$ with
maximum property satisfies all six lower-rung assertions:
\begin{enumerate}
\item $V(T')$ is elementary under its generating coloring;
\item under every proper coloring leaving exactly $e_0$ uncolored and
stable on its final ladder closure, that closure is interchangeable for
two colors whenever at least one of them is missing there;
\item stability descends unconditionally at each PE support to the prefix before the
support and, when its first-appearance index is greater than one, also to
the preceding ladder stage; stored pairs at
distinct PE slots with distinct supports and the same first-appearance
index are disjoint;
\item at a terminal PE, the stored two-color component through the
support meets the terminal ladder tree only at that support;
\item every nonminimum endpoint of a defective boundary color precedes
each later PE support or SE extension vertex allowed by its
first-appearance index; and
\item every proper terminal stable coloring leaving exactly $e_0$
uncolored is a mod-coloring of the ETT.
\end{enumerate}

Let $\sigma_n$ be proper, leave exactly $e_0$ uncolored, and be
$T_n\oplus R_n,D_n,\varphi_n$-stable.

For arbitrary $\alpha,\beta\in[k]$, if
$\alpha\in\overline\sigma_n(T_n\vee R_n)$ and $\{\alpha\}$ is
closed on $T_n\vee R_n$ under $\sigma_n$, then $\alpha,\beta$ are
$T_n\vee R_n$-interchangeable under $\sigma_n$.
All boundaries, trees, and two-color components use the original edge
type, so distinct parallel edges retain their identities.
\end{lemma}
\begin{proof}
Write $X=V(T_n)$, $Y=V(R_n)$, and
$U=V(T_n\vee R_n)$, with vertex sets understood as in
\noderef{TreeSequenceVertexSet}.  The generating-color field of
\noderef{ExtendedTashkinovTree} identifies $\varphi_n$ with the last
coloring of the certified series.  Eliminate the last certificate in
\noderef{TashkinovSeries} using the terminal-PE hypothesis.  The PE
constructor of \noderef{TashkinovSeriesStep} supplies the exact support
$v_n$, exact connecting-edge identity $f_n$, and colors
$\gamma_n,\delta_n$.  It also supplies properness and exact
one-uncoloredness of $\varphi_n$, and the equations recording the
terminal $\gamma_n,\delta_n$-switch.  No endpoint pair is substituted
for $f_n$.

The parallel clause of \noderef{HierarchyBase} says that $R_n$ is a
$\varphi_n$-closure of the ordered prefix $T_n(v_n)$, and supplies
an \noderef{OrderedTreeUnion} certificate with
\[
                         U=X\cup Y.                       \tag{1}
\]
It also says that the missing-color set of $T_n$ under the base
coloring is its missing-color set under $\varphi_n$.
The two conjuncts of \noderef{CombinedStableColoring}, with stability
expanded as in \noderef{StableColoring}, say that
$\sigma_n$ is $(T_n,D_n,\varphi_n)$-stable and simultaneously
$(R_n,\varnothing,\varphi_n)$-stable.  Applying
\noderef{CombinedStableClosurePreserved} gives, for every displayed edge
identity $g$ of either tree,
\[
             \sigma_n(g)=\varphi_n(g),                    \tag{2}
\]
and proves that the same ordered tree $R_n$, with the same edge
identities, is a $\sigma_n$-closure of $T_n(v_n)$.

We shall need the delta-boundary assertion under a stable coloring rather
than merely under $\varphi_n$, so we establish that transport now.  Write
$\pi$ for the maximum-defective coloring and $\rho$ for the second stable
coloring in the terminal PE constructor of
\noderef{TashkinovSeriesStep}.  The choice certificate
\noderef{MaximumDefectiveChoice} gives an exact edge $f_n$ incident with
$v_n$ such that $\pi(f_n)=\delta_n$.  It also puts $v_n$ in $X$.  The rule
\noderef{ParallelExtensionColor} says that $\gamma_n$ is missing at
$v_n$ under $\pi$.  The vertex clause of the stable relation from $\pi$
to $\rho$ in \noderef{StableColoring} therefore says that $\gamma_n$ is
missing at $v_n$ under $\rho$.  Since $f_n$ is incident with $T_n$ by
\noderef{TreeIncidentEdge}, and $\delta_n$ belongs to the enlarged
control set of that stable relation, its edge clause gives
\[
                       \rho(f_n)=\delta_n.                    \tag{2a}
\]
Thus $\gamma_n\ne\delta_n$.  There is no $\gamma_n$-edge incident with
$v_n$, by the complement definition \noderef{MissingColors}.  Moreover,
$f_n$ is the only $\delta_n$-edge incident with $v_n$: any different such
edge would be adjacent to $f_n$ by \noderef{AdjacentEdges} and would
contradict properness in \noderef{IsProperEdgeColoring}.

The edge-node of $f_n$ is adjacent to the vertex-node $v_n$ in
\noderef{TwoColorIncidenceGraph}, so it lies in the full
$(\gamma_n,\delta_n)$-component of $v_n$ in
\noderef{TwoColorComponent}.  The terminal-color equation says that
$\varphi_n$ is the \noderef{TwoColorComponentSwap} of $\rho$ on exactly
that full component.  It changes $f_n$ from $\delta_n$ to $\gamma_n$.
It cannot change any other edge incident with $v_n$ to $\delta_n$, since
such an edge would have had color $\gamma_n$ under $\rho$, and there is no
such edge; every incident edge carrying neither switch color is fixed.
Consequently
\[
              \delta_n\in\overline\varphi_n(v_n).             \tag{2b}
\]

More generally, let $\tau$ be any proper coloring leaving exactly $e_0$
uncolored which is $(T_n,D_n,\varphi_n)$-stable.  Its vertex clause at
$v_n$ transfers (2b) to
$\delta_n\in\overline\tau(v_n)$.  The witness $v_n\in X$ and the union
definition \noderef{TreeMissingColors} then give
$\delta_n\in\overline\tau(T_n)$.  The PE constructor supplies properness
and exact one-uncoloredness of $\varphi_n$, while these are assumptions on
$\tau$.  Package the two colorings with these two properties and apply the
symmetry clause of \noderef{sc2}; it says that $\varphi_n$ is
$(T_n,D_n,\tau)$-stable.

Now fix an exact identity
$g\in\partial_{\tau,\delta_n}(T_n)$.  Unfolding the color filter
\noderef{BoundaryColorEdges} retains both
$g\in\partial_G(X)$ and $\tau(g)=\delta_n$.  By
\noderef{BoundaryEdges}, the first membership supplies an endpoint
$p\in X$ and the other endpoint outside $X$.  The endpoint relation is
\noderef{Incident}, and hence the inside endpoint witnesses
\noderef{TreeIncidentEdge} for $g$ and $T_n$.  Because $\delta_n$ is in
$\overline\tau(T_n)$, the edge clause of the reverse stable relation gives
\[
                 \varphi_n(g)=\tau(g)=\delta_n.                \tag{2c}
\]
The geometric boundary membership is unchanged, and a second use of
\noderef{BoundaryColorEdges} therefore gives
\[
 \partial_{\tau,\delta_n}(T_n)
       \subseteq\partial_{\varphi_n,\delta_n}(T_n).             \tag{2d}
\]
This derivation applies both to the original $\sigma_n$ and to either
outside-swap coloring used below, once its properness,
one-uncoloredness, and forward stability have been checked.

We record three preliminary facts, with their hypotheses checked at the
coloring at which they will be used.

First, $U$ is elementary under $\sigma_n$.  Indeed, retain the same
core, closure, and support position from the hierarchy base.  Equality
of missing-color sets on $T_n$ is the vertex clause of the first stable
relation.  The preceding paragraph supplies the required
$\sigma_n$-closure, and the old ordered-union certificate is unchanged.
These data form a hierarchy base under $\sigma_n$.  The first stable
relation, properness, exact one-uncoloredness, positivity, maximum
property, the terminal-PE equation, and the displayed six lower-rung
assertions are precisely the hypotheses of
\noderef{ParallelClosureUnionElementary}.  Hence
\[
                  U\text{ is elementary under }\sigma_n.   \tag{3}
\]

By \noderef{InterchangePECommonMissingClosed}, every color
$c\in\overline\sigma_n(T_n)\cap\overline\sigma_n(R_n)$ is closed
on $U$:
\[
 c\in\overline\sigma_n(T_n)\cap\overline\sigma_n(R_n)
 \Longrightarrow \partial_{\sigma_n,c}(U)=\varnothing.       \tag{5}
\]
Precisely, that node returns the singleton instance of
\noderef{TreeColorSetClosed}; evaluating its universal clause at the sole
member $c\in\{c\}$ and unfolding \noderef{BoundaryColorEdges} gives the
displayed equality of exact boundary-edge sets.  Conversely, whenever we
use membership $c\in\overline\tau(S)$ below, the definition
\noderef{TreeMissingColors} supplies a displayed vertex $x\in V(S)$ with
$c\in\overline\tau(x)$.  We will name those witnesses rather than treating
tree-missing membership as vertexwise missingness without explanation.
By \noderef{InterchangePEClosureConstituent},
\[
 a\in\overline\sigma_n(R_n)
 \Longrightarrow a,b\text{ are $R_n$-interchangeable}.        \tag{7}
\]
The last PE certificate already extracted above identifies exact
$v_n,\gamma_n,\delta_n$ and gives the last-slot chosen-vertex, gamma,
and delta equalities.  Applying
\noderef{InterchangePETreeAwayFromDelta} with those equalities yields
\[
 a\in\overline\sigma_n(T_n),\quad a\ne\delta_n,\quad
 b\ne\delta_n
 \Longrightarrow a,b\text{ are $T_n$-interchangeable}.        \tag{8}
\]

We prove (iii) by contradiction, beginning with the component bookkeeping
which will also be used below.  Expanding the negation of
\noderef{TreeInterchangeable}, suppose $a$ is missing on $U$, is closed
there, and there are roots $z_1,z_2$ whose full incidence components in
\noderef{TwoColorComponent}
$P_1,P_2$ are unequal, each meets $U$, and each satisfies
\noderef{TwoColorPathComponent}.  Unpacking either path condition gives
$a\ne b$, and each component has a node with at most one neighbor.  By
\noderef{TwoColorComponentsDisjoint}, the unequal full components are
node-disjoint.

Here the hypotheses of the present lemma say, literally, that
$a\in\overline\sigma_n(U)$ and that $U$ is
$\{a\}$-closed in the sense of \noderef{TreeColorSetClosed}.
The former has a vertex witness by \noderef{TreeMissingColors}; the latter,
at the singleton member $a$, says
$\partial_{\sigma_n,a}(U)=\varnothing$, with this exact colored boundary
set defined by \noderef{BoundaryColorEdges}.  Every later assertion that
an $a$-edge cannot cross the boundary uses this equality.

For every $P_i$ which contains an edge-node, its path certificate and
that witness give \noderef{TwoColorNontrivialPathComponent}.  Apply
\noderef{TwoColorComponentSpanningPath} to $z_i$; its simple spanning path
has support exactly $P_i$.  We verify the nature of its two ends.  By
\noderef{TwoColorIncidenceGraph}, every neighbor of an edge-node is a
vertex-node incident with that exact edge identity.  The two ends of an
edge are distinct by \noderef{LooplessMultigraph}, and incidence is the
membership relation of \noderef{Incident}; hence every colored edge-node
has exactly two distinct vertex-node neighbors.  On the other hand,
properness gives the degree bound of
\noderef{TwoColorIncidenceDegreeAtMostTwo} at every incidence node.  The
at-most-one-neighbor node supplied by
\noderef{TwoColorPathComponent} lies in a nontrivial connected component,
so it has exactly one neighbor and is a vertex-node.  Consequently the
spanning path is not closed by an extra adjacency between its two ends.
Any other extra neighbor of a path end would already occur in the spanning
support and would give either that internal node or the path end degree at
least three.  Thus both spanning-path ends have incidence degree one.  They
are vertex-nodes, since edge-nodes have degree two.

At such a vertex end $x$, at least one of $a,b$ is missing.  Indeed, if
both were present, \noderef{PresentColors} would give incident exact edge
identities of the two respective colors.  They are distinct because
$a\ne b$, and \noderef{TwoColorIncidenceGraph} would give $x$ two
neighbors, contrary to its degree one.  The two ends of every nontrivial
$P_i$ are therefore distinct vertex-nodes missing $a$ or $b$.

Only now use (3), with elementariness unfolded by
\noderef{ColoredVertexSetElementary}.  If both $P_i$ are nontrivial,
their four ends are distinct by component disjointness, and at least two
lie outside $U$: otherwise three tree ends would assign the same one of
$a,b$ to two distinct vertices of $U$.  If, after relabeling, $P_1$ has
no edge-node, its connectedness and its meeting with $U$ make it the
singleton component at some $x\in U$.  Both $a$ and $b$ are missing at
$x$, since an incident edge of either color would be an adjacent edge-node
in $P_1$.  Elementariness makes $P_2$ nontrivial and forces both ends of
its spanning path outside $U$.  The two components cannot both be
singletons, because their distinct tree vertices would repeat both
missing colors.  Thus in all cases there are at least two outside ends of
nontrivial component paths.  From each selected end, follow its spanning
path up to its first vertex in $U$; this exists because each $P_i$ meets
$U$.  If the ends belong to different components, the resulting exit
subpaths are disjoint by the preceding component disjointness.  If they
are the two ends of one component, take the two terminal portions up to
the first and last encounters with $U$.  Their inside ends are distinct:
otherwise the two distinct boundary identities at that common vertex
would both have color $b$, since the closed color $a$ has no boundary
identity, contradicting properness.  Hence these portions too are
vertex-disjoint.  Denote the resulting exit subpaths by $Q_1,Q_2$ and
their distinct outside vertex ends by $r_1,r_2$.  Each $r_i$ is one of the
nontrivial full-component ends just analyzed.

We may reduce to
\[
        a\in\overline\sigma_n(T_n)
             \cap\overline\sigma_n(R_n).                     \tag{9}
\]
Here is the reduction with the recoloring hypotheses explicit.  The
support $v_n$ lies in both constituents: it lies in the ordered prefix
and that prefix lies in its closure.  Since $k\geq\Delta(G)+1$ by
\noderef{KTriple}, the number of colors exceeds the number of exact edge
identities incident with $v_n$: the latter number is at most
$\Delta(G)$ by \noderef{MaximumDegree}.  Properness permits at most one
incident edge of each present color, so the complement definition
\noderef{MissingColors}, together with the edge-witness definition
\noderef{PresentColors}, gives a color $c$ missing at $v_n$.
As $v_n\in X\cap Y$, the same vertex is a witness in both unions of
\noderef{TreeMissingColors}; hence $c$ is common missing.  Formula (5),
expanded through \noderef{TreeColorSetClosed} and
\noderef{BoundaryColorEdges}, says exactly that neither a $c$-colored nor
an $a$-colored exact edge identity lies in the boundary of $U$.
Apply the outside-tree transposition
\noderef{OutsideTreeColorSwap} to $(a,c)$: it fixes every exact edge
identity incident with $U$ and interchanges $(a,c)$ on every other edge
identity.  We verify its coloring properties here because only the two
displayed colors, rather than every color missing on $U$, are assumed
closed.  There is no $a$- or $c$-colored boundary identity by the two
exact empty-boundary equalities just obtained.  At a vertex of $U$ all incident
values are fixed.  At a vertex outside $U$, all nonboundary incident
values undergo the same transposition, while every boundary value avoids
$(a,c)$ and is fixed.  Thus injectivity of the transposition and the
adjacent-edge clause of \noderef{IsProperEdgeColoring} prove properness
of the new coloring.  The transposition fixes the uncolored value, so
exactly $e_0$ remains uncolored by \noderef{IsOneEdgeUncolored}.  It also
fixes all incident values and hence all missing-color
sets on $U$, by the complement and witness definitions
\noderef{MissingColors} and \noderef{PresentColors}.  Its edge values are
equal on every identity incident with either constituent, so both clauses
of \noderef{StableColoring} hold there for every control set.

We first show that $b$ is present at $v_n$.  If it were missing there,
then $v_n\in X\cap Y$ would, by \noderef{TreeMissingColors}, put $b$ in
both constituent tree-missing sets.  Equation (5), unfolded through
\noderef{TreeColorSetClosed} and \noderef{BoundaryColorEdges}, would give
$\partial_{\sigma_n,b}(U)=\varnothing$.  We already have
$\partial_{\sigma_n,a}(U)=\varnothing$.  A connected
$(a,b)$-component which meets $U$ could then have no vertex outside $U$:
the first exact edge identity on an incidence walk from an inside vertex
to an outside vertex belongs geometrically to \noderef{BoundaryEdges},
and its color filter in \noderef{BoundaryColorEdges} would make it an
$a$- or $b$-colored boundary identity.  Thus both components would be trapped in
$U$.  If both are nontrivial, their four distinct ends lie in $U$ and
repeat one of the two missing colors.  If one is a singleton, its tree
vertex misses both colors, while an end of the other component is a
distinct tree vertex missing one of them.  Either possibility contradicts
(3) by \noderef{ColoredVertexSetElementary}.
So $b$ is present at $v_n$, whereas $c$ is missing there; in particular
$c\ne b$.

On each $Q_i$, the first geometric boundary identity belongs to the set
in \noderef{BoundaryEdges}.  It is fixed and has color $b$, because its
alternating color is $a$ or $b$ and the exact set
$\partial_{\sigma_n,a}(U)$ supplied by \noderef{TreeColorSetClosed} and
\noderef{BoundaryColorEdges} is empty.  Every later vertex lies outside $U$, so every
later edge identity of this exit portion is outside $U$ and its alternating
$(a,b)$-colors become alternating $(c,b)$-colors.  At the outside end
$r_i$, the new $(c,b)$-neighbor set is exactly the old $(a,b)$-neighbor
set: nonincident identities undergo the transposition, while identities
incident with $U$ have neither color $a$ nor $c$ by the two empty exact
colored-boundary sets and are fixed.  Hence $r_i$ still has degree one and belongs
to the full $(c,b)$-component containing the transformed $Q_i$.

Let $K$ be the full $(c,b)$-component through $v_n$.  Since $c$ is
missing and $b$ is present there, properness gives $v_n$ exactly one
$(c,b)$-neighbor.  Thus $K$ contains an exact $b$-edge and, by
\noderef{TwoColorNontrivialPathComponent}, is a nontrivial path component.  It can
contain at most one of $Q_1,Q_2$.  Otherwise it contains the three
distinct degree-one vertex-nodes $v_n,r_1,r_2$, whereas its spanning path
from \noderef{TwoColorComponentSpanningPath} has only two ends.  Retain a
transformed $Q_i$ not contained in $K$, and let $L$ be its full component.
The exact edge-node on $Q_i$ and its degree-one outside end $r_i$ verify
that $L$ also satisfies \noderef{TwoColorNontrivialPathComponent}.  It meets $U$ at
the inside end of $Q_i$.  The dichotomy in
\noderef{TwoColorComponentsDisjoint} makes $K$ and $L$ disjoint, since
otherwise they would be equal and $Q_i$ would lie in $K$.  We have
therefore produced two distinct full nontrivial $(c,b)$-path components
meeting $U$, while $c$ is missing at the common vertex
$v_n\in X\cap Y$.  Replacing the coloring and $a$ by this new coloring
and $c$ yields (9).  The stable relation with the original coloring composes with the two
constituent stable relations by \noderef{sc2}.  Thus the replacement
coloring is proper, leaves exactly $e_0$ uncolored, and is forward stable
from $\varphi_n$ on $T_n$ and $R_n$.  In particular, the general
calculation (2b)--(2d) applies to it.  It again satisfies the hypotheses of
\noderef{InterchangePECommonMissingClosed},
\noderef{InterchangePEClosureConstituent}, and
\noderef{InterchangePETreeAwayFromDelta}; reapplying those three nodes
gives (5), (7), and (8) for the replacement.  Equation (2), the
hierarchy base under the replacement, and
\noderef{ParallelClosureUnionElementary} likewise give (3).  We
henceforth rename this coloring $\sigma_n$ and rename the two
constructed components $P_1,P_2$.

By \noderef{TreeMissingColors}, choose a witness in $X$ and a witness in
$Y$ for the two memberships in (9).  By (3) and
\noderef{ColoredVertexSetElementary}, those witnesses must be the same
vertex; otherwise $U$ would have two vertices missing $a$.  Call the
common witness $v$.
Take $P_1$ to be the component through $v$.  By (7), every other
$(a,b)$-component meeting $R_n$ equals $P_1$, so $P_2$ is disjoint
from $R_n$.  It nevertheless meets $T_n$.  Thus $(a,b)$ are not
$T_n$-interchangeable, and (8) implies
\[
                         a=\delta_n\quad\text{or}\quad
                         b=\delta_n.                           \tag{10}
\]

Suppose first that $b=\delta_n$.  The delta-boundary clause of
\noderef{HkuParallelBoundaryStructure} is stated for $\varphi_n$, so we
first perform the exact transport proved in (2b)--(2d).  Namely,
$\delta_n$ is missing at $v_n$ under $\varphi_n$ by (2b), and the vertex
clause of the present forward stable relation makes it missing at $v_n$
under the current $\sigma_n$.  Since $v_n\in X$,
\noderef{TreeMissingColors} gives
$\delta_n\in\overline\sigma_n(T_n)$.  Properness and
\noderef{IsOneEdgeUncolored} hold for both $\varphi_n$ and $\sigma_n$;
symmetry in \noderef{sc2} therefore supplies the reverse
$(T_n,D_n,\sigma_n)$-stable relation.

Take an arbitrary exact
$g\in\partial_{\sigma_n,\delta_n}(T_n)$.  By
\noderef{BoundaryColorEdges}, $g$ retains its membership in the geometric
set \noderef{BoundaryEdges} and satisfies $\sigma_n(g)=\delta_n$.
The geometric membership gives an inside endpoint $p\in X$; by
\noderef{Incident} this endpoint witnesses
\noderef{TreeIncidentEdge}.  The reverse edge clause of
\noderef{StableColoring} now gives
$\varphi_n(g)=\sigma_n(g)=\delta_n$.  Thus the same exact $g$ lies in
$\partial_{\varphi_n,\delta_n}(T_n)$.  Only now apply
\noderef{HkuParallelBoundaryStructure}: it supplies a tree vertex $w$
incident with $g$ and strictly preceding $v_n$.  The positional definition
\noderef{TreeVertexOrder}, together with the prefix definition
\noderef{TreeSequencePrefix}, puts $w$ in the ordered prefix $T_n(v_n)$.
The \noderef{HierarchyBase} closure certificate, expanded through
\noderef{TashkinovClosure}, \noderef{TashkinovAugmentationSequence},
\noderef{TashkinovAugmentation}, and
\noderef{TreeSequenceInitialSegment}, retains every prefix vertex in
$R_n$.  Hence every current $\delta_n$-boundary identity of $T_n$ is
incident with a vertex of $Y$.
Fix one of the nontrivial $a,\delta_n$-components meeting $T_n$.  If it
has a vertex outside $T_n$, a shortest incidence walk from a vertex in
$T_n$ to one outside it uses an $a$- or $\delta_n$-boundary identity, by
the exact geometric/color filters in \noderef{BoundaryEdges} and
\noderef{BoundaryColorEdges}.
An $a$-boundary identity cannot leave $U$, because
$\partial_{\sigma_n,a}(U)=\varnothing$ by
\noderef{TreeColorSetClosed};
if it leaves $T_n$, its outside endpoint is therefore in $U-X\subseteq
Y$ by (1).  A $\delta_n$-boundary identity has an inside endpoint in the
prefix and hence in $Y$, by the preceding paragraph.  In either subcase
the component meets $R_n$.

It remains possible that every vertex-node of the component lies in $X$.
Apply the endpoint analysis above to this component.  Its two distinct
vertex ends each miss $a$ or $\delta_n$.  The common witness $v$ from
(9) lies in $X\cap Y$ and misses $a$.  The explicit terminal calculation
(2a)--(2b) proves that $\delta_n$ is missing at $v_n$ under $\varphi_n$;
the vertex clause of \noderef{StableColoring} transfers this to the
present $\sigma_n$.  The support belongs to the ordered prefix and hence
$v_n\in X\cap Y$.  By elementariness (3), unfolded through
\noderef{ColoredVertexSetElementary}, an end which misses $a$ must
equal $v$, and an end which misses $\delta_n$ must equal $v_n$.  Thus
every end belongs to $\{v,v_n\}\subseteq Y$ (or the distinctness of the
ends already makes this internal case impossible).  In either event the
component cannot be disjoint from $R_n$.  Consequently every
nontrivial component meeting $T_n$ meets $R_n$.  In particular both
$P_1,P_2$ meet $R_n$, contradicting (7).

Suppose instead that $a=\delta_n$.  Equation (5) says that
$\delta_n$ is closed on $U$; by \noderef{TreeColorSetClosed} and
\noderef{BoundaryColorEdges}, this means the exact set
$\partial_{\sigma_n,\delta_n}(U)$ is empty.

We next construct the second common missing color without suppressing the
root-edge argument.  Put
$x_0=T_n.\operatorname{vertex}(0)$ and
$x_1=T_n.\operatorname{vertex}(1)$.  These indices exist because every
\noderef{TreeSequence} has positive length.  Its first-edge clause says
$T_n.\operatorname{edge}(0)=e_0$.  Its attaches clause at index zero has
an earlier index in the one-element set, hence that index is zero and
\[
                  \operatorname{ends}(e_0)=\{x_1,x_0\}.
\]
Its vertex-injectivity clause gives $x_0\ne x_1$.  Thus these are the two
distinct exact endpoints of $e_0$, up to the immaterial order in the
unordered endpoint pair.  The image definition
\noderef{TreeSequenceVertexSet} puts both in $X$.

Let $r$ be the support position recorded by \noderef{HierarchyBase}.
Its support-positive clause gives $1\leq r$, so the prefix definition
\noderef{TreeSequencePrefix} retains the vertex entries at indices zero
and one.  The same parallel clause makes $R_n$ a
\noderef{TashkinovClosure} of this prefix.  Expanding that definition gives
a \noderef{TashkinovAugmentationSequence}; each one-step
\noderef{TashkinovAugmentation} retains its old tree as a
\noderef{TreeSequenceInitialSegment}.  Induction along the finite sequence
therefore retains $x_0$ and $x_1$ in its final tree $R_n$.  Hence
\[
                         x_0,x_1\in X\cap Y.                  \tag{11}
\]

We count missing colors at $x_0$.  By the endpoint equation above,
$e_0$ is incident with $x_0$.  Let $I$ be the finite set of exact edge
identities incident with $x_0$.  The definition \noderef{Degree} gives
$|I|=d_G(x_0)$, and \noderef{MaximumDegree} gives
$d_G(x_0)\leq\Delta(G)$.  The definition
\noderef{IsOneEdgeUncolored} says that $e_0$ has value uncolored and every
other member of $I$ has some color.  Every color in
\noderef{PresentColors} at $x_0$ is supplied by such a member of
$I-\{e_0\}$.  Conversely, each member of $I-\{e_0\}$ supplies a present
color, and two distinct members supply distinct colors: they are adjacent
at $x_0$ by \noderef{AdjacentEdges}, so equality of their colors would
contradict \noderef{IsProperEdgeColoring}.  Thus
\[
 |\operatorname{PresentColors}(x_0)|=|I|-1=d_G(x_0)-1.
\]
In particular
$|\operatorname{PresentColors}(x_0)|+2\leq\Delta(G)+1\leq k$, where the
last inequality is the numerical clause of \noderef{KTriple}.  It is
available here because certification at slot zero in
\noderef{TashkinovSeries} forces the initial constructor of
\noderef{TashkinovSeriesStep}, which supplies that $k$-triple.  Since
\noderef{MissingColors} is the complement of the present-color set in the
$k$ colors, at least two colors are missing at $x_0$.  A set with at least
two elements cannot be contained in the singleton $\{\delta_n\}$, so
choose a missing color $c\ne\delta_n$.  By (11), the same $x_0$ is a
vertex witness in both union definitions of \noderef{TreeMissingColors};
hence
\[
 c\in\overline\sigma_n(T_n)\cap\overline\sigma_n(R_n),
 \qquad c\ne\delta_n.
\]
By (5), evaluated through \noderef{TreeColorSetClosed} and
\noderef{BoundaryColorEdges}, the exact $c$-colored boundary of $U$ also
is empty.  Apply
\noderef{OutsideTreeColorSwap} to $(\delta_n,c)$.  The color $\delta_n$ is
tree-missing by (9), the color $c$ is missing at the common vertex $x_0$,
and both exact colored boundaries are empty.  Thus the same edge-by-edge
mixed-adjacency check used above proves that the result is proper and
leaves exactly $e_0$ uncolored.  It fixes every edge identity incident
with $U$, hence \noderef{PresentColors} and \noderef{MissingColors} give
the same missing-color set at every vertex of $U$; both clauses of
\noderef{StableColoring} hold on each constituent relative to the old
coloring.  Transitivity in \noderef{sc2} composes these relations with the
old forward relations from $\varphi_n$.  Therefore the new coloring again
satisfies the stable hypotheses of
\noderef{InterchangePECommonMissingClosed},
\noderef{InterchangePEClosureConstituent}, and
\noderef{InterchangePETreeAwayFromDelta}; these give (5), (7), and (8)
for the new coloring.  Since all vertex missing sets on $U$ are unchanged,
(3) also holds for it.  We now regard this new coloring as the current
coloring in the remainder of the argument.

We repeat, at this new swap, the full-component check rather than merely
reusing vertex-disjointness.  Apply the endpoint construction above afresh
to the current $P_1,P_2$, obtaining two vertex-disjoint exit portions
$Q'_1,Q'_2$ with distinct outside degree-one ends $r'_1,r'_2$.  By (3)
and \noderef{ColoredVertexSetElementary},
any witnesses that the chosen $c$ is missing on $T_n$ and on $R_n$ must
coincide.  The explicit common witness is still $q=x_0$, because the
outside swap fixes every identity incident with $U$.  If $b$ were missing
at $q$, then \noderef{TreeMissingColors} would make $b$ common missing,
and (5), via \noderef{TreeColorSetClosed} and
\noderef{BoundaryColorEdges}, would give
$\partial_{\sigma_n,b}(U)=\varnothing$.  The colors $\delta_n$ and $b$
would both be missing and have empty exact colored boundaries on $U$, so
the \noderef{BoundaryEdges} and four-end argument
given above would put all four distinct ends of $P_1,P_2$ in $U$ and
contradict \noderef{ColoredVertexSetElementary}.  Thus $b$ is present at
$q$, $c$ is missing
there, and $c\ne b$.

The $c,b$-component $K'$ through $q$ is therefore nontrivial and has
$q$ as a degree-one vertex, so it satisfies
\noderef{TwoColorNontrivialPathComponent}.  Each $Q'_i$ becomes a $(c,b)$-exit
portion: its boundary identity has color $b$, all later
$\delta_n$-values become $c$, and its outside end remains degree one by
the same exact-neighbor comparison.  The spanning-path argument forbids
$K'$ from containing both $r'_1,r'_2$.  Choose a transformed $Q'_i$ not
in $K'$.  Its full component $L'$ contains its exact edge-node, has the
degree-one end $r'_i$, and meets $U$, hence is another
\noderef{TwoColorNontrivialPathComponent}.  By
\noderef{TwoColorComponentsDisjoint}, $K'$ and $L'$ are distinct full
components.  The root component $K'$ meets both $X$ and $Y$ at $q$, while
$L'$ meets $U=X\cup Y$ at the inside end of its retained exit portion.  If
that inside end lies in $X$, then $K',L'$ contradict (8), because
$c\ne\delta_n$ and $b\ne\delta_n$ in this branch.  If it lies in $Y$,
then they contradict (7), since $c$ is missing on $R_n$.  Both
alternatives in (10) are impossible, proving (iii).
\end{proof}
